\documentclass[12pt,a4paper]{article}
\usepackage[utf8]{inputenc}
\usepackage[T1]{fontenc}
\usepackage{amsmath,amssymb,amsthm}
\usepackage{physics}
\usepackage{geometry}
\usepackage{enumitem}
\usepackage{hyperref}
\usepackage{fancyhdr}
\usepackage{titlesec}

\geometry{margin=1in}
\pagestyle{fancy}
\fancyhf{}
\rhead{Lecture 1}
\lhead{Advanced Quantum Mechanics}
\cfoot{\thepage}

\titleformat{\section}{\large\bfseries}{}{0em}{}
\titleformat{\subsection}{\normalsize\bfseries}{}{0em}{}

\newtheorem{definition}{Definition}

\title{\textbf{Lecture 1: Review of Quantum Mechanics} \\ \large Symmetries, the Schr\"odinger Equation, and Generators}
\date{}

\begin{document}
\maketitle
\thispagestyle{fancy}

% ============================================================
\section{1. Review of the Schr\"odinger Equation}
% ============================================================

The time-dependent Schr\"odinger equation governs how the state vector $\ket{\psi(t)}$ evolves:
\[
    i\hbar \frac{\partial}{\partial t}\ket{\psi(t)} = H \ket{\psi(t)},
\]
where $H$ is the Hamiltonian operator.  The formal solution is
\[
    \ket{\psi(t)} = e^{-iHt/\hbar}\ket{\psi(0)}.
\]
The operator $U(t) = e^{-iHt/\hbar}$ is unitary, ensuring that the norm $\braket{\psi}{\psi} = 1$ is preserved at all times.

\subsection{1.1 Time-independent Schr\"odinger equation}
If we expand $\ket{\psi}$ in energy eigenstates $H\ket{E_n} = E_n\ket{E_n}$, then the time evolution becomes:
\[
    \ket{\psi(t)} = \sum_n c_n \, e^{-iE_n t/\hbar}\ket{E_n},
\]
where $c_n = \braket{E_n}{\psi(0)}$.  Each component oscillates at its own frequency $\omega_n = E_n/\hbar$.

\subsection{1.2 Expectation values}
For any observable $L$, its expectation value evolves as:
\[
    \frac{d}{dt}\expval{L} = \frac{i}{\hbar}\expval{[H, L]}.
\]
If $[H, L] = 0$, then $\expval{L}$ is a constant of the motion---the observable is \emph{conserved}.

% ============================================================
\section{2. Symmetry in Quantum Mechanics}
% ============================================================

\begin{definition}[Symmetry]
A \textbf{symmetry} of a quantum system is a unitary operator $U$ that commutes with the Hamiltonian:
\[
    [U, H] = 0.
\]
Equivalently, if $\ket{\psi(t)}$ is a solution of the Schr\"odinger equation, then $U\ket{\psi(t)}$ is also a solution.
\end{definition}

Symmetries come in two varieties:
\begin{itemize}[nosep]
    \item \textbf{Discrete symmetries}: reflections, parity, permutations.  The operator $U$ acts by a finite transformation.
    \item \textbf{Continuous symmetries}: translations, rotations.  The operator $U(\epsilon)$ depends on a continuous parameter and can be built up from infinitesimal steps.
\end{itemize}

% ============================================================
\section{3. Continuous Symmetries and Generators}
% ============================================================

A continuous symmetry is generated by a Hermitian operator.  For an infinitesimal transformation parameterized by $\epsilon$:
\[
    U(\epsilon) = I - i\epsilon G,
\]
where $G = G^\dagger$ is the \textbf{generator} of the symmetry.  The condition $[U, H] = 0$ implies $[G, H] = 0$, so the generator is a conserved quantity.

\subsection{3.1 Translation symmetry and momentum}

Translation by a distance $\delta$ acts on wave functions as:
\[
    T(\delta)\psi(x) = \psi(x - \delta).
\]
Expanding to first order:
\[
    \psi(x - \delta) \approx \psi(x) - \delta \frac{\partial \psi}{\partial x} = \left(1 - \delta \frac{\partial}{\partial x}\right)\psi(x).
\]
Comparing with $U(\delta) = I - i\delta G/\hbar$, we identify the generator:
\[
    G = -i\hbar\frac{\partial}{\partial x} = p,
\]
which is the \textbf{momentum operator}.  If the Hamiltonian is translationally invariant ($[T(\delta), H] = 0$), then momentum is conserved.

\subsection{3.2 Noether's theorem (quantum version)}

The connection between symmetries and conservation laws:
\begin{enumerate}[nosep]
    \item Every continuous symmetry has an associated Hermitian generator $G$.
    \item $[U, H] = 0 \;\Longleftrightarrow\; [G, H] = 0$.
    \item $[G, H] = 0$ means $G$ is conserved: $\frac{d}{dt}\expval{G} = 0$.
\end{enumerate}
This is the quantum-mechanical version of Noether's theorem: \textbf{every continuous symmetry implies a conservation law}, and vice versa.

% ============================================================
\section{4. Finite Transformations from Infinitesimal Ones}
% ============================================================

A finite translation by $a$ is built by iterating infinitesimal translations:
\[
    T(a) = \lim_{N\to\infty}\left(1 - \frac{ia}{\hbar N}p\right)^N = e^{-iap/\hbar}.
\]
More generally, for a symmetry generated by $G$ with parameter $\theta$:
\[
    U(\theta) = e^{-i\theta G/\hbar}.
\]
This is a unitary operator because $G$ is Hermitian.

% ============================================================
\section{5. Summary of Key Correspondences}
% ============================================================

\begin{center}
\begin{tabular}{l l l}
\textbf{Symmetry} & \textbf{Generator} & \textbf{Conserved Quantity} \\
\hline
Time translation & Hamiltonian $H$ & Energy \\
Space translation & Momentum $p$ & Momentum \\
Rotation & Angular momentum $L$ & Angular momentum \\
\end{tabular}
\end{center}

\medskip
\noindent The unifying principle: the Hamiltonian generates time evolution, and any operator that commutes with it is conserved.  The structure of symmetries and their generators is the foundation on which the rest of the course is built.

\end{document}
